\section{Representation, evaluation order, and provenance}
\label{sec:effects-provenance}
The same relation can have several representations with different observable
operations. Consequently, an equation between relational expressions is only
one component of a compiler replacement argument. This section isolates that
component, gives a finite account of the distinction witnessed by
Counterexample~\ref{cx:effects}, and develops the provenance question raised by
embedded computations. The finite model and its propositions are mathematical
analysis of the existing routes. The sorted representation and provenance
instrumentation discussed below are not implemented MiniLinq features.

\subsection{An abstraction function does not specify every observation}
\label{sec:representation-observations}
For a concrete storage representation $r$, let $\mathcal S_r$ be its valid
states and let
\[
  a_r:\mathcal S_r\longrightarrow \operatorname{Rel}(U)
\]
map a valid state to its represented relation of sort $U$. A representation
invariant states when $s\in\mathcal S_r$ and how its operations implement that
relation. For example, the HashSet invariant includes component equality and
hashing laws; a sorted-vector invariant additionally requires a specified
ordering. Equal abstractions $a_r(s)=a_{r'}(s')$ do not identify the observable
behaviors of the two states.

Let $\Omega_r(s)$ be a nonempty set of permitted complete scan sequences for
each valid state $s$. A scan of
set storage must satisfy
\[
 \bar t\in\Omega_r(s)\quad\Longrightarrow\quad
 \set(\bar t)=a_r(s)\ \land\ |\bar t|=|a_r(s)|.
\]
Here tuple values are compared after representation abstraction. A particular
scan selects one permitted sequence and consumes its successive prefixes.
There is no fresh arbitrary choice of the next row at every poll. An
unspecified-order API admits an abstract contract containing all permutations;
this is an overapproximation of implementations, not a claim that one fixed
HashSet can actually realize every permutation.

\begin{center}\small
\begin{tabular}{>{\raggedright\arraybackslash}p{.19\linewidth}>{\raggedright\arraybackslash}p{.31\linewidth}>{\raggedright\arraybackslash}p{.42\linewidth}}
\toprule
Semantic boundary & Required invariant & Additional observable commitment \\
\midrule
Named RA result & A finite relation of its declared sort & No enumeration or
callback schedule is selected by the relational denotation. \\
HashSet storage & Unique abstract tuples; lawful equality and hashing & A scan
visits each tuple once, in an order the API does not prescribe. \\
Sorted vector (comparison) & The same unique tuples, increasing under a total
order consistent with equality & A full scan has that order. A key-only order
still permits choices within equal-key groups. \\
Access or Pipeline cursor & Bindings, residual iterator state, and the specified
local control flow & A demand determines which sources and callbacks are
evaluated before the next answer or exhaustion. \\
\bottomrule
\end{tabular}
\end{center}
The HashSet API supplies the relevant unspecified-order contract~\cite{hashset}.
The sorted row is a contrasting representation contract, not a new language
level. Whether an invariant is recorded in an IR, established by a pass, or
maintained by a concrete library is a separate implementation question.

For a finite successful client protocol $d$, write
$\mathcal B_d(P,I,\sigma)$ for the set of behaviors admitted by program $P$,
the input representation contract, and initial host store $\sigma$.
The choices in a behavior include scan sequences consistent with its storage
and complete successive-poll history. Let $O$ select observations from a
behavior. Three useful choices are:
\begin{enumerate}
\item $O_{\rm set}$: the set of answers, for a client that exhausts the cursor;
\item $O_{\rm seq}$: the complete answer sequence, again for an exhaustive client;
\item $O_{\rm client}$: the answers and host events at each build, query, poll,
and drop request, together with observable store changes.
\end{enumerate}
These choices are explicit projections of a behavior, not competing meanings
of the word relation. Failed and infinite behaviors require the additional
semantics in C-Effects; the notation here is restricted to finite successful
requests.

For the finite refinement claims below we require \emph{response totality}:
for each admitted input and client, $\mathcal B_d(P,I,\sigma)\ne\emptyset$
for both programs, and every admitted maximal execution responds successfully
to each request in the claimed finite client protocols. This is a separate
premise, not a consequence of retaining only successful behaviors. Without
it an always-panicking target has an empty successful-behavior set and would
satisfy inclusion vacuously. A semantics admitting failure and divergence
must instead account for them explicitly; C-Effects remains open.

For an admitted-input predicate $\mathsf{Pre}$, a target $T$ refines a source
$S$ at observation $O$ when, for every $(I,\sigma)$ satisfying $\mathsf{Pre}$
and every client $d$ in the claimed client class,
\[
 \{O(b)\mid b\in\mathcal B_d(T,I,\sigma)\}
 \subseteq
 \{O(b)\mid b\in\mathcal B_d(S,I,\sigma)\}.
\]
Equivalence requires equality. Thus a sorted scan can refine an abstract scan
that permits every permutation, provided construction and polling satisfy the
same remaining contract. This statement does not license changing callback
timing, representation-visible values, or the set of demanded computations.

The relational theorems use $O_{\rm set}$ and $\mathcal H_{\rm rel}$: all
admitted successful complete behaviors then have the same observed answer set.
An Access-to-Pipeline theorem can state a stronger poll-level property because
that edge retains a particular operational structure. Different edges therefore
have different source invariants and preservation properties. A composition
argument must state which observations survive each edge; it cannot silently
compose set preservation with an unrestricted contextual-equivalence claim.

\paragraph{Sorted output does not determine preceding effects.}
The existing \code{materialize} methods sort the answers after query evaluation.
This fixes the order of the returned vector under its ordering assumptions.
It does not impose an order on earlier scans, earlier calls to \code{let}, or
earlier guards. In Counterexample~\ref{cx:effects}, sorting two lazy answers and
one materialized answer leaves different sets. Choosing a sorted intermediate
store could fix a scan order, but still would not equate eager and interleaved
evaluation over that order.

\subsection{A finite state-and-trace account of clause evaluation}
\label{sec:effect-sequencing}
The essential distinction is between a value and the computation that produces
it, as made explicit by Moggi's computational semantics~\cite{prov-moggi}.
To analyze the example without claiming a semantics for all Rust, fix a set
of stores $\mathcal W$ and a set of observable event labels $E$. Consider
total, deterministic callbacks in
\[
 \mathcal M A=\mathcal W\longrightarrow A\times\mathcal W\times E^*.
\]
Fixing an admitted scan sequence resolves the enumeration choice for this
finite model. The ambient Rust environment $\rho$ is fixed in these equations;
$\nu$ supplies its logical-variable bindings. Callback evaluation has no panic, divergence, concurrency,
allocation identity, or destructor effects. Those exclusions delimit this
model, not the syntax accepted by the compiler.

The unit and sequencing operations are
\[
\begin{aligned}
 \operatorname{return}(v)(\sigma)&=(v,\sigma,\epsilon),\\
 (m\mathbin{\gg=} f)(\sigma)&=(w,\sigma_2,\alpha\concat\beta),\\
 &\quad\text{where }m(\sigma)=(v,\sigma_1,\alpha),\quad
 f(v)(\sigma_1)=(w,\sigma_2,\beta).
\end{aligned}
\]
Associativity permits regrouping successive computations; it does not permit
exchanging them. An expression with free logical variables in $X$ denotes,
in this model, a function from valuations over $X$ to a computation
$\mathcal M A$, rather than necessarily a function from valuations to $A$.

A clause $c$ maps one current valuation to a computation producing a finite
sequence of successor valuations. For an irrefutable pattern $p$ binding fresh
variables, let $\operatorname{match}_p(v)$ be its binding valuation. Define
\[
\begin{aligned}
 T_{\mathsf{let}\ p=e}(\nu)
   &= H_e(\nu)\mathbin{\gg=}
       (\lambda v.\operatorname{return}[\nu\cup\operatorname{match}_p(v)]),\\
 T_{\mathsf{if}\ e}(\nu)
   &= H_e(\nu)\mathbin{\gg=}
       (\lambda b.\operatorname{return}(\mathsf{if}\ b\
          \mathsf{then}\ [\nu]\ \mathsf{else}\ [])),\\
 T_A(\nu)&=\operatorname{return}(\operatorname{matches}_A(\nu)).
\end{aligned}
\]
The last line is restricted to a relational atom whose membership tests are
silent; its fixed enumeration lists compatible extensions. Computed atom
arguments require their own expression evaluation sites before the associated
test. These three clauses suffice for the finite witnesses below; they do not
assign missing semantics to negation or aggregation.

For $f:A\to\mathcal M[B^*]$, define ordered traversal with concatenation:
\[
\begin{aligned}
 \operatorname{trav}(f,[]) &=\operatorname{return}[],\\
 \operatorname{trav}(f,a::\bar a)
  &= f(a)\mathbin{\gg=}(\lambda\bar b.\
       \operatorname{trav}(f,\bar a)\mathbin{\gg=}
       (\lambda\bar c.\operatorname{return}(\bar b\concat\bar c))).
\end{aligned}
\]
For a clause list $C$, the two completed-evaluation schedules are:
\[
\begin{aligned}
 \operatorname{depth}([],\nu)&=\operatorname{return}[\nu],\\
 \operatorname{depth}(c::C,\nu)
  &=T_c(\nu)\mathbin{\gg=}
       (\lambda\bar\mu.\operatorname{trav}
           (\lambda\mu.\operatorname{depth}(C,\mu),\bar\mu)),\\[2pt]
 \operatorname{bulk}([],\bar\nu)&=\operatorname{return}\bar\nu,\\
 \operatorname{bulk}(c::C,\bar\nu)
  &=\operatorname{trav}(T_c,\bar\nu)\mathbin{\gg=}
       (\lambda\bar\mu.\operatorname{bulk}(C,\bar\mu)).
\end{aligned}
\]
Depth completes the suffix for one successor before moving to the next. Bulk
completes one clause for every prefix before starting the next clause. Bulk
here retains sequence order and duplicates, so even this comparison isolates
interleaving from HashSet normalization. The materialized implementation also
collects sets and selects a permitted scan of each result. In the witness all
extended rows are distinct, so those extra choices do not remove the difference.
The cursor rules in Section~\ref{sec:cursor} refine the depth schedule to partial
demand; these complete-evaluation equations alone do not specify when build or
the first poll performs a computation.

\paragraph{The counterexample as an execution derivation.}
Fix scan order $10,20$, let $T_x$ denote \code{tick(x)}, and let $G_x$
denote \code{latest(n)} using the value produced for that row. The counter
is initially zero. Each table entry records a state transition and a result.
\begin{center}\small
\begin{tabular}{lll}\toprule
Schedule & Callback sequence & Surviving $x$ values \\
\midrule
Depth & $T_{10}:0\mapsto1;\ G_{10}:1=1;\
          T_{20}:1\mapsto2;\ G_{20}:2=2$ & $10,20$ \\
Bulk & $T_{10}:0\mapsto1;\ T_{20}:1\mapsto2;\
         G_{10}:2\ne1;\ G_{20}:2=2$ & $20$ \\
\bottomrule
\end{tabular}
\end{center}
Reversing the input order makes the other materialized row survive. Changing
the scan order of the extended store cannot make both survive: the store is
already at two when either guard runs. The two schedules are thus
distinguishable for every admitted ordering in the example. This proves the
absence of unrestricted route equivalence. It does not by itself choose the
source language's intended effect schedule or identify which route violates
that as-yet-unspecified contract.

\subsection{Which equalities an optimization actually needs}
\label{sec:effect-laws}
Free-variable availability answers whether a callback can be supplied its
logical arguments. It does not answer whether moving, merging, or suppressing
its evaluation is valid. Even preserving the written order between clause
\emph{sites} leaves the interleaving between their \emph{dynamic invocations}
undetermined.

\begin{definition}[Independent finite computations]
Two computations $m:\mathcal M A$ and $n:\mathcal M B$ are independent at
the claimed reachable stores if
\[
\begin{aligned}
 &m\mathbin{\gg=}(\lambda a.\
 n\mathbin{\gg=}(\lambda b.\operatorname{return}(a,b)))\\
 ={}&n\mathbin{\gg=}(\lambda b.\
 m\mathbin{\gg=}(\lambda a.\operatorname{return}(a,b))).
\end{aligned}
\]
Equality includes both returned components, the final store, and the event
sequence. Neither computation's input is supplied by the other's result.
\end{definition}

A useful sufficient condition in a finite-cell store model is as follows.
Suppose a callback reads only a footprint $R$, writes only a footprint $W$,
and its result, writes, and events depend only on its explicit arguments and
the initial values at $R$. The footprint guarantees hold for all reachable
stores under comparison, not just the accesses observed in one run. For
callbacks $m,n$, require
\[
 W_m\cap(R_n\cup W_n)=\emptyset,
 \qquad W_n\cap(R_m\cup W_m)=\emptyset.
\]
Assume no other effects, and that the two event sequences commute under
concatenation (silence suffices). Each callback then reads unchanged inputs
when exchanged; their disjoint updates commute; their individual results and
combined event sequence agree. This proves the independence equation in this
restricted model. Global variables, aliasing, opaque calls, and macro bodies
prevent deriving these premises from logical free variables alone.

\begin{proposition}[Finite two-clause interchange]
\label{prop:interchange}
Fix $n$ prefix valuations and an order-preserving materialization. For each
prefix $i$, let $E_i$ be a total extension and $G_i$ its following total guard.
No other operation changes the prefix population. Suppose $G_i$ and $E_j$
are independent for every $i<j$ at every reachable store of the exchanges.
Then bulk evaluation $E_1\cdots E_nG_1\cdots G_n$ and depth evaluation
$E_1G_1\cdots E_nG_n$ have the same complete answer sequence, final store,
and event sequence.
\end{proposition}
\begin{proof}
Move $G_1$ left over $E_n,\ldots,E_2$. Each adjacent exchange preserves
the two labeled results, store, and events by independence, and leaves $G_1$
after the $E_1$ that supplies its argument. Then move $G_2$ left over
$E_n,\ldots,E_3$, and continue. The finite process reaches the depth schedule.
Guard results and prefix order are preserved, so filtering keeps the same
sequence. No evaluation is added or removed.
\end{proof}
This proposition deliberately says nothing about build timing or stopping
after the first answer. It also requires order-preserving materialization,
which a HashSet does not guarantee. It provides a local condition an effectful
fusion argument could use, rather than asserting that the existing routes
satisfy it. Saarikivi et al.'s effectful comprehensions provide an established
example of proving fusion for an explicitly modeled class of stateful stream
computations~\cite{prov-fusion}; accepting unrestricted Rust supplies no
corresponding model automatically.

\paragraph{Projection introduces a different obligation.}
Let $R(x,y)=\{(1,10),(1,20)\}$ and consider
\begin{lstlisting}
answer(x) :- R(x,y), let ignored:() = notify(x);
\end{lstlisting}
Let notify append its argument to a visible log and return $()$. Evaluating
the let once for each valuation of $\{x,y\}$ appends $1,1$. Projecting to $x$
and eliminating duplicates before the let appends $1$. The argument's value
is the same, and $y$ is absent from both the head and the expression's free
variables, but the evaluation populations differ. This is a semantic example
of a proposed rewrite, not a claim that this particular rewrite occurs in the
current compiler. It shows why relational liveness alone does not establish
effect-preserving column pruning.

Deleting an unused let requires that discarding its computation preserve the
selected observations. Merging two identical evaluations requires a sharing
law. Moving an evaluation across a guard can change whether it occurs at all;
moving it from poll to build can change whether a client demands it. These
are separate obligations from exchange. In particular, a let whose pattern
binds no variables is still a computation; its empty binding set does not
prove it erasable. For the pure total fragment the required laws follow from
$\mathcal H_{\rm rel}$, whereas arbitrary effectful expressions need a more
specific admissibility argument.

\subsection{What provenance must distinguish}
\label{sec:provenance-model}
Three existing lines of work answer different questions. Semiring provenance
annotates positive relational derivations using addition for alternatives
and multiplication for joint use of inputs~\cite{prov-semiring}. Dependency
provenance asks how changes to input components can affect output components,
and provides sound approximations rather than generally computable minimal
dependencies~\cite{prov-dependency}. Provenance traces record operational
derivations and support statements about what actually happened and how a
trace adapts when inputs change~\cite{prov-traces}. The present problem is
closest to the latter two when host computations can inspect shared state.

For example, tag two matching alternatives with input tokens $r_1s_1$ and
$r_2s_2$. The polynomial $r_1s_1+r_2s_2$ records their relational derivations.
It does not specify which match caused the first host call, whether two
extensions were completed before either guard, or whether a client requested
the second answer. Its commutative operations intentionally support
relational equivalences. Merely changing multiplication to a noncommutative
operator does not supply stores, dynamic invocation counts, stage boundaries,
or a semantics for evaluating a guard. Those objects must be defined.

For a finite sequential execution in the model above, an explanatory record
can be constructed without changing the query IR. Let
\[
 \mathcal P=(N,\lambda,D,C,F,\prec).
\]
Here $N$ contains input-tuple occurrences, generated valuation occurrences,
callback entry/return events, guard decisions, yields, and client requests. Each
dynamic occurrence has a fresh identifier: two occurrences can carry equal
valuations without being the same occurrence. The label $\lambda$ records
the source clause site, invocation identity, relevant input and result
values, and the client stage. The edge relations have different meanings:
\begin{description}[style=nextline,leftmargin=1.5em]
\item[$D$: data dependence.]
A generated value has an incoming edge from its producer or input occurrence;
an invocation has incoming edges from the producers of its supplied logical
arguments. For an opaque callback this is a conservative argument dependence,
not a claim that every supplied argument semantically influences its result.
\item[$C$: control dependence.]
A suffix invocation or yield has an edge from the guard decision or iteration
occurrence that enabled it. A rejected guard is retained even though it yields
no answer. This is dynamic enabling in the defined execution, not an assertion
that every such edge is a minimal counterfactual cause.
\item[$F$: state dependence.]
In an instrumented finite-cell host semantics, each read records its location,
value, and the latest preceding write to that location, or its initial value.
This gives a read-from edge even when two callbacks share no logical variable.
Repeated writes also create successive state versions. A callback with several
accesses is expanded into its ordered access events; it is not falsely treated
as one atomic memory operation.
\item[$\prec$: recorded execution order.]
The strict total creation order on all nodes in $N$ restricts to the actual
event order, including build/query/poll boundaries, callback invocation, and
store accesses. Input and initial-store nodes are placed first in an
arbitrary fixed order; generated value nodes are created before their
consumers. This order contains scheduling facts as
well as dependence constraints. An earlier event is not thereby a cause of
every later event. Reordering additionally needs to account for write-after-read
and write-after-write conflicts, even where no value flows along a read-from
edge.
\end{description}
The directions of $D,C,F$ are from predecessor to consumer. Input values and
the initial store precede events. All edges therefore respect $\prec$.
The graph is finite and acyclic for a finite execution, although its source
program or query hypergraph need not be acyclic.

This record is a mathematical extraction from an instrumented execution, not
an implemented provenance collector. Exact host reads are not available from
Syn syntax or the current free-variable helper. Treating an unknown host
computation as depending on the whole observable store gives a conservative
analysis boundary within the finite model; it supplies no missing
instrumentation or replay theorem.

\paragraph{The distinguishing edge in the running counterexample.}
In the depth execution, $T_{10}$ produces both $n=1$ and the counter version
read by $G_{10}$. In the bulk execution, $T_{10}$ still produces the argument
$n=1$, but $T_{20}$ produces the counter version read by $G_{10}$. Thus the
guard's explicit data dependence is unchanged while its state dependence
changes. In the following display only, an invocation-level arrow abbreviates
the corresponding producer/consumer path through result or access events:
\[
\begin{array}{lll}
 \text{depth:}&T_{10}\xrightarrow{D}G_{10},&T_{10}\xrightarrow{F}G_{10},\\
 \text{bulk:}&T_{10}\xrightarrow{D}G_{10},&T_{20}\xrightarrow{F}G_{10}.
\end{array}
\]
The second-row invocation affects whether the first row is returned. Input
tuple lineage for the first row alone cannot show this dependency. A
per-binding invocation identity also explains the notify example: projecting
and deduplicating before evaluation merges two valuation occurrences into one,
and consequently deletes one invocation and one event. Finally, retaining
client-request nodes explains why an unconsumed lazy suffix has no invocation
at all, even though its input tuple exists.

The state edges have a simple constructive definition. Maintain a map
$v:\mathsf{Loc}\to N$ from each finite host-store location to its current
version node. Initially create one node labeled with $\sigma(\ell)$ for each
location $\ell$. On a primitive read or write, append a fresh event identifier
after all previous events and apply, respectively,
\[
\begin{aligned}
 \mathsf{read}(\ell):\quad
 &\lambda(r)=\mathsf{read}(\ell,\sigma(\ell)),\quad
 F'=F\cup\{(v(\ell),r)\},\quad v'=v;\\
 \mathsf{write}(\ell,a):\quad
 &\lambda(w)=\mathsf{write}(\ell,a),\quad
 \sigma'=\sigma[\ell\mapsto a],\quad v'=v[\ell\mapsto w].
\end{aligned}
\]
A read leaves the store unchanged; a write leaves $F$ unchanged. Data edges
record the producers of the written value. The separate recorded execution
order retains conflicts between successive accesses even where $F$ has no
edge. At callback entry and return, add its argument and result edges; at a
guard or iteration decision, add the corresponding enabling edges for the
events it starts. All other components are retained. This is instrumentation
of the finite primitive host semantics, which must actually expose its reads
and writes to construct the record.

To state erasure, let $z$ be an underlying machine state containing the
residual computation, valuation, store and client state. Its instrumented
state is $(z,\mathcal P,v)$, with erasure
$\operatorname{erase}(z,\mathcal P,v)=z$. Each underlying transition is
decorated once by appending its finite event/value nodes in order and updating
the version map as above. Instrumentation never determines an underlying
branch or changes a host value. This construction presupposes the primitive
host transitions expose the required access and enabling information.

\begin{proposition}[Finite record consistency]
Decorate a finite successful run by retaining every event in order, the
values returned by that event, its enabling predecessors, and the actual
read-from information of the finite-cell host semantics. Projecting the
instrumented states by erase gives the same underlying run; every recorded read observes the latest
preceding version of its location; and every recorded invocation has the
arguments and control decisions supplied by that run.
\end{proposition}
\begin{proof}
Induct on the event sequence. Initially the recorded versions are the initial
store. A read uses the current version by the host transition; a write
establishes the next version. Argument and control edges refer to values or
decisions already produced. Appending their labels changes no execution
state, and removing those labels recovers the original transition. The
induction establishes all three properties, including the final observations.
\end{proof}
This is a consistency property of faithful instrumentation of one finite run.
Erasure projects an augmented execution; the record $\mathcal P$ alone is
not asserted to reconstruct the residual machine or replay the program.
It is weaker than a theorem permitting replay after changing inputs: a changed
guard can create invocations absent from the original record. Provenance-trace
work makes precisely such fidelity obligations explicit~\cite{prov-traces}.
Nor does graph reachability by itself establish an account of actual causes;
that requires a specified intervention semantics~\cite{prov-causality}.

\subsection{The resulting research and compiler obligations}
The examples motivate finer observations of dynamic query evaluation. They do
not establish a previously unknown kind of provenance. The concrete research
question is whether an observation-indexed provenance semantics can expose
exactly the distinctions each existing IR makes, while permitting justified
abstractions at compilation boundaries. A candidate development must discharge
the following obligations:
\begin{enumerate}
\item Define the source effect contract, including enumeration freedom,
per-valuation evaluation population, and demand boundaries. This selects the
observations against which a route is judged.
\item Define each provenance view and its erasure. Relational annotations,
invocation histories, and state-dependency explanations must have stated
semantics; one does not follow merely by naming graph edges.
\item Prove instrumentation adequacy and a separate preservation theorem for
each lowering. A pass that duplicates or fuses invocations needs a relation
between source and target records, not a presumed bijection of their nodes.
\item Prove any claimed dependency soundness, changed-input fidelity, or
intervention property over its explicit host fragment. Finite-run consistency
alone supplies none of these stronger guarantees.
\end{enumerate}
The present manuscript establishes the local finite distinctions and conditional
interchange result above. It leaves the full provenance construction,
all-Rust effect contract, instrumentation, and cross-IR provenance preservation
open. These are substantive semantic tasks; neither a parser extension nor an
extra administrative label discharges them.
